\documentclass[10pt,conference,compsoc]{IEEEtran}

\usepackage{graphicx}
\usepackage[nocompress]{cite}
\usepackage{amsmath}

\begin{document}

\title{Covariate-Independent Individuality and Subject Re-Identification in CGM
Postprandial Glucose Excursions}

\author{\IEEEauthorblockN{Sun Jung}
\IEEEauthorblockA{International School Ho Chi Minh City\\
High School Student\\
sun.jung.co2027@gmail.com}}

\maketitle

\begin{abstract}
Continuous glucose monitors record how a person responds to every meal, but it is
unclear how much of that response is specific to the person rather than to the
food. Using the public CGMacros dataset, we describe each of 1{,}657 meal
excursions from 45 subjects with ten features covering amplitude, timing, and
shape, and treat each meal as a repeated measurement of its subject. After
adjusting for demographics and fasting labs, peak and baseline glucose remain the
most person-specific features, meals from the same person cluster together more
tightly than meals from different people (permutation $p = 0.003$), and a
nearest-neighbour model re-identifies the subject behind a held-out meal well
above chance. The same features predict diabetes under subject-grouped
cross-validation, and the features that are most individual are also the most
predictive. Postprandial glucose responses therefore carry a covariate-independent
individual signature that doubles as a signal of metabolic health.
\end{abstract}

\begin{IEEEkeywords}
continuous glucose monitoring, individuality, re-identification, intraclass
correlation, diabetes classification
\end{IEEEkeywords}

\section{Introduction}

Continuous glucose monitors are now cheap enough that people without diabetes
wear them, and each sensor produces a dense record of how the body handles food.
The most informative part of that record is the meal response: glucose rises
after eating, reaches a peak, and returns toward baseline. The height, timing,
and steepness of that curve differ from person to person, and this paper asks a
simple question about those differences. Is the shape of a meal response
individual to the person, in a way that ordinary explanations like age and body
weight cannot account for, and if so, does that same individual signal say
anything about their metabolic health?

Earlier work has shown that people mount different responses to the same food and
that continuous glucose patterns fall into distinct groups. What it has largely
left untested is whether the individual component survives once demographics and
laboratory values are removed, and it has not tried to measure that individuality
the way biometrics does, by asking whether an unseen meal can be traced back to
the right person. Framing the problem as re-identification is what lets us report
individuality as an accuracy rather than an impression.

This paper makes two contributions. First, using covariate-adjusted
repeatability, multivariate distance, and a re-identification test on held-out
meals, we show that postprandial excursion profiles carry an individual signal
that demographics and labs do not explain. Second, we show that the same features
that make a person identifiable also predict diabetes under subject-grouped
cross-validation, so individuality and clinical signal turn out to live in the
same features rather than in separate ones.

\section{Related Work}

Continuous glucose monitoring records interstitial glucose every few minutes and
has changed how day-to-day metabolism is observed~\cite{danne2017}. In routine
use, though, most of that signal is compressed into a handful of static
summaries. Time in range and the glucose management indicator report how often
glucose sits inside a target band or estimate an A1c-like average~%
\cite{battelino2019,bergenstal2018}, and variability indices add a measure of
spread~\cite{monnier2008}. These numbers are useful for management, but they
mostly describe where glucose is rather than how it moves. A meal response has a
shape: a rise, a peak, and a return toward baseline, and that shape is largely
lost once a day is reduced to an average. Clinical guidelines have argued for
years that the postprandial period itself carries information about metabolic
health~\cite{ceriello2008}, which suggests the excursion is worth modeling on its
own terms.

That the same meal can produce very different responses in different people is by
now well documented. Zeevi et al.~\cite{zeevi2015} found that postprandial
responses to identical foods vary widely across individuals and can be predicted
from personal and microbiome features, which motivated personalized rather than
uniform dietary advice. Hall et al.~\cite{hall2018} took the idea further,
clustering CGM traces into distinct ``glucotypes'' that separated people by how
variable their glucose was, including people without a diabetes diagnosis. Both
results establish that individual differences are real. What they leave open is
whether those differences are stable and specific enough to act as a signature of
the person once the obvious explanations, such as age, body mass, and blood
chemistry, are accounted for.

Measuring how much of a trait belongs to the person is an old problem. The
intraclass correlation coefficient (ICC) splits the variance of repeated
measurements into a between-subject part and a within-subject part, so a high ICC
means a feature is consistent within a person yet differs across people. A
related but stricter view comes from biometrics, where the question is not only
whether people differ on average but whether an unseen sample can be matched back
to the right individual. Posing CGM meal responses this way, as a
re-identification task, turns a loose notion of individuality into a measurable
quantity: how often the correct person appears among the top guesses for a
held-out meal. To our knowledge this framing has not been applied to postprandial
glucose curves, even though recognizing a person from a physiological signal is a
familiar goal elsewhere.

CGM-derived features have also been used to separate diabetic from non-diabetic
individuals, which is consistent with the underlying physiology: impaired insulin
secretion and insulin resistance push peaks higher and slow the return to
baseline after eating~\cite{defronzo2004}. One caution runs through this line of
work. When the diabetes label is itself defined by fasting glucose or A1c,
features that essentially re-measure glucose level will track that label almost by
construction. The more informative question is whether the shape and timing of the
excursion, and not only its height, add anything beyond the diagnostic
thresholds. We keep this distinction in mind when we later separate amplitude
features from shape and kinetic ones.

\end{document}
